\chapter{2003年上海卷（秋理）}
\section{填空题}
\begin{problem}
函数 \( y = \sin x \cos \left( x + \frac{\pi}{4} \right) + \cos x \sin \left( x + \frac{\pi}{4} \right) \) 的最小正周期 \( T = \)\fillinblank{}.
\end{problem}

\begin{answer}
\(\pi\)
\end{answer}

\begin{problem}
若 \(x = \frac{\pi}{3}\) 是方程 \(2\cos (x + \alpha) = 1\) 的解，其中 \(\alpha \in \left(0, 2\pi\right)\)，则 \(\alpha =\)\fillinblank{}.
\end{problem}

\begin{answer}
\(\frac{4\pi}{3}\)
\end{answer}

\begin{problem}
在等差数列 \(\{a_{n}\}\) 中，\(a_{5} = 3\)，\(a_{6} = -2\)，则 \(a_{4} + a_{5} + \cdots + a_{10} =\fillinblank{}\).
\end{problem}

\begin{answer}
\(-49\)
\end{answer}

\begin{problem}
在极坐标系中，定点 \(A\left(1, \frac{\pi}{2}\right)\)，点 \(B\) 在直线 \(\rho \cos \theta + \rho \sin \theta = 0\) 上运动，当线段 \(AB\) 最短时，点 \(B\) 的极坐标是\fillinblank{}.
\end{problem}

\begin{answer}
\(\left(\frac{\sqrt{2}}{2}, \frac{3\pi}{4}\right)\)
\end{answer}

\begin{problem}
在正四棱锥 \(P - ABCD\) 中，若侧面与底面所成二面角的大小为 \(60^{\circ}\)，则异面直线 \(PA\) 与 \(BC\) 所成角的大小等于\fillinblank{}. （结果用反三角函数值表示）
\end{problem}

\begin{answer}
\(\arctan 2\)
\end{answer}

\begin{problem}
设集合 \(A = \{x \mid |x| < 4\}\)，\(B = \{x \mid x^2 - 4x + 3 > 0\}\)，则集合 \(\{x \mid x \in A \text{且} x \notin A \cap B\} =\fillinblank{}\).
\end{problem}

\begin{answer}
\([1,3]\)
\end{answer}

\begin{problem}
在 \(\triangle ABC\) 中，\(\sin A: \sin B: \sin C = 2:3:4\)，则 \(\angle ABC =\)\fillinblank{}. （结果用反三角函数值表示）
\end{problem}

\begin{answer}
\(\arccos\frac{11}{16}\)
\end{answer}

\begin{problem}
若首项为 \(a_1\)，公比为 \(q\) 的等比数列 \(\{a_n\}\) 的前 \(n\) 项和总小于这个数列的各项和，则首项 \(a_1\)，公比 \(q\) 的一组取值可以是 \((a_1, q) =\fillinblank{}\).
\end{problem}

\begin{answer}
\(\left(1,\frac{1}{2}\right)\)
\end{answer}

\begin{problem}
某国际科研合作项目成员由 \(11\) 个美国人，\(4\) 个法国人和 \(5\) 个中国人组成，现从中随机选出两位作为成果发布人，则此两人不属于同一个国家的概率为\fillinblank{}. （结果用分数表示）
\end{problem}

\begin{answer}
\(\frac{119}{190}\)
\end{answer}

\begin{problem}
方程 \(x^{3} + \lg x = 18\) 的根 \(x \approx\fillinblank{}\). （结果精确到 \(0.1\)）
\end{problem}

\begin{answer}
\(2.6\)
\end{answer}

\begin{problem}
已知点 \(A\left(0, \frac{2}{n}\right)\)，\(B\left(0, -\frac{2}{n}\right)\)，\(C\left(4 + \frac{2}{n}, 0\right)\)，其中 \(n\) 为正整数. 设 \(S_{n}\) 表示 \(\triangle ABC\) 外接圆的面积，则 \(\displaystyle\lim_{n \to \infty} S_{n} = \)\fillinblank{}.
\end{problem}

\begin{answer}
\(4\pi\)
\end{answer}

\begin{problem}
给出问题： \(F_{1}\)， \(F_{2}\) 是双曲线 \( \frac{x^{2}}{16}-\frac{y^{2}}{20}=1 \) 的焦点，点 \(P\) 在双曲线上. 若点 \(P\) 到焦点 \(F_{1}\) 的距离等于 \(9\)，求点 \(P\) 到焦点 \(F_{2}\) 的距离. 某学生的解答如下：双曲线的实轴长为 \(8\)，由 \( \left|PF_{1}\right| - \left|PF_{2}\right| = 8 \)，即 \( \left|9 - \left|PF_{2}\right|\right| = 8 \)，得 \( \left|PF_{2}\right| = 1 \) 或 \(17\). 该学生的解答是否正确？若正确，请将他的解题依据填在下面空格内，若不正确，将正确的结果填在下面空格内\fillinblank{}.
\end{problem}

\begin{answer}
\(17\)
\end{answer}

\section{单选题}
\begin{problem}
下列函数中，既为偶函数又在 \((0, \pi)\) 上单调递增的是（\quad）

\choices
    {\( y = \tan |x| \)}
    {\( y = \cos (-x) \)}
    {\( y = \sin \left(x - \frac{\pi}{2}\right) \)}
    {\( y = \left|\cot \frac{x}{2}\right| \)}
\end{problem}

\begin{answer}
\(C\)
\end{answer}

\begin{problem}
在下列条件中，可判断平面 \(\alpha\) 与 \(\beta\) 平行的是（\quad）

\choices
    {\(\alpha\)，\(\beta\) 都垂直于平面 \(\gamma\)}
    {\(\alpha\) 内存在不共线的三点到 \(\beta\) 的距离相等}
    {\(l\)，\(m\) 是 \(\alpha\) 内两条直线，且 \(l \parallel \beta\)，\(m \parallel \beta\)}
    {\(l\)，\(m\) 是两条异面直线，且 \(l \parallel \alpha\)，\(m \parallel \alpha\)，\(l \parallel \beta\)，\(m \parallel \beta\)}
\end{problem}

\begin{answer}
\(D\)
\end{answer}

\begin{problem}
\(a_{1}\)，\(b_{1}\)，\(c_{1}\)，\(a_{2}\)，\(b_{2}\)，\(c_{2}\) 均为非零实数，不等式 \(a_{1}x^{2} + b_{1}x + c_{1} > 0\) 和 \(a_{2}x^{2} + b_{2}x + c_{2} > 0\) 的解集分别为集 \(M\) 和 \(N\)，那么 \(\frac{a_1}{a_2} = \frac{b_1}{b_2} = \frac{c_1}{c_2}\) 是“\(M = N\)”的（\quad）

\choices
    {充分非必要条件}
    {必要非充分条件}
    {充要条件}
    {既非充分又非必要条件}
\end{problem}

\begin{answer}
\(D\)
\end{answer}

\begin{problem}
\(f(x)\) 是定义在区间 \([-c, c]\) 上的奇函数，其图象如图所示：令 \(g(x) = a f(x) + b\)，则下列关于函数 \(g(x)\) 的叙述正确的是（\quad）

\bitmapfigure[max width=0.5\linewidth]{img/2003/shanghai_science/q16_fig1.png}

\choices
    {若 \(a < 0\)，则函数 \(g(x)\) 的图象关于原点对称}
    {若 \(a = -1\)，\(-2 < b < 0\)，则方程 \(g(x) = 0\) 有大于 2 的实根}
    {若 \(a \neq 0\)，\(b = 2\)，则方程 \(g(x) = 0\) 有两个实根}
    {若 \(a \geqslant 1\)，\(b < 2\)，则方程 \(g(x) = 0\) 有三个实根}
\end{problem}

\begin{answer}
\(B\)
\end{answer}

\section{解答题}
\begin{problem}
已知复数 \(z_{1} = \cos \theta - \i\)，\(z_{2} = \sin \theta + \i\)，求 \(|z_{1} \cdot z_{2}|\) 的最大值和最小值.
\end{problem}

\begin{solution}
令 \(c=\cos\theta\)，\(s=\sin\theta\)，则 \(c^{2}+s^{2}=1\)，且 \(z_{1}z_{2}=(c-\i)(s+\i)=1+cs+\i(c-s)\). 因此
\[
|z_{1}z_{2}|^{2}=(1+cs)^{2}+(c-s)^{2}=2+c^{2}s^{2}=2+\frac{1}{4}\sin^{2}2\theta.
\]
由于 \(0\leqslant\sin^{2}2\theta\leqslant1\)，所以 \(2\leqslant|z_{1}z_{2}|^{2}\leqslant\frac{9}{4}\)，从而 \(|z_{1}z_{2}|\) 的最小值为 \(\sqrt{2}\)，最大值为 \(\frac{3}{2}\). 当 \(\sin2\theta=0\) 时取到最小值，当 \(|\sin2\theta|=1\) 时取到最大值.
\end{solution}

\begin{problem}
已知平行六面体 \(ABCD - A_{1}B_{1}C_{1}D_{1}\) 中，\(A_{1}A \perp\) 平面 \(ABCD\)，\(AB = 4\)，\(AD = 2\). 若 \(B_{1}D \perp BC\)，直线 \(B_{1}D\) 与平面 \(ABCD\) 所成的角等于 \(30^{\circ}\)，求平行六面体 \(ABCD - A_{1}B_{1}C_{1}D_{1}\) 的体积.

\bitmapfigure[max width=0.58\linewidth]{img/2003/shanghai_science/q18_fig1.png}
\end{problem}

\begin{solution}
设 \(\overrightarrow{AB}=\bs{b}\)，\(\overrightarrow{AD}=\bs{d}\)，\(\overrightarrow{AA_{1}}=\bs{u}\). 由于 \(A_{1}A\perp\) 平面 \(ABCD\)，有 \(\bs{u}\cdot\bs{d}=0\). 又 \(\overrightarrow{B_{1}D}=\bs{d}-\bs{b}-\bs{u}\)，\(\overrightarrow{BC}=\bs{d}\)，由 \(B_{1}D\perp BC\) 得 \(0=(\bs{d}-\bs{b}-\bs{u})\cdot\bs{d}=\bs{d}^{2}-\bs{b}\cdot\bs{d}\)，所以 \(\bs{b}\cdot\bs{d}=|\bs{d}|^{2}=4\). 平行四边形 \(ABCD\) 的面积为 \(\sqrt{|\bs{b}|^{2}|\bs{d}|^{2}-(\bs{b}\cdot\bs{d})^{2}}=\sqrt{4^{2}\cdot2^{2}-4^{2}}=4\sqrt{3}\).

直线 \(B_{1}D\) 在底面 \(ABCD\) 上的射影为 \(\overrightarrow{BD}=\bs{d}-\bs{b}\)，故其长度满足 \(|\bs{d}-\bs{b}|=\sqrt{4+16-2\times4}=2\sqrt{3}\). 设 \(AA_{1}=h\)，由直线与平面所成角为 \(30^{\circ}\)，有 \(\tan30^{\circ}=h/(2\sqrt{3})\)，所以 \(h=2\). 因而平行六面体的体积为 \(4\sqrt{3}\times2=8\sqrt{3}\).
\end{solution}

\begin{problem}
已知数列 \(\{a_{n}\}\) （\(n\) 为正整数）是首项是 \(a_{1}\)，公比为 \(q\) 的等比数列.

\begin{enumerate}
\item 求和：\(a_{1}\mathrm C_{2}^{0} - a_{2}\mathrm C_{2}^{1} + a_{3}\mathrm C_{2}^{2}\)，\(a_{1}\mathrm C_{3}^{0} - a_{2}\mathrm C_{3}^{1} + a_{3}\mathrm C_{3}^{2} - a_{4}\mathrm C_{3}^{3}\).
\item 由前一问的结果归纳概括出关于正整数 \(n\) 的一个结论，并加以证明；
\end{enumerate}
\end{problem}

\begin{solution}
\begin{enumerate}
\item 由 \(a_{2}=a_{1}q\)，\(a_{3}=a_{1}q^{2}\)，\(a_{4}=a_{1}q^{3}\)，第一个和为 \(a_{1}-2a_{2}+a_{3}=a_{1}-2a_{1}q+a_{1}q^{2}=a_{1}(1-q)^{2}\)，第二个和为 \(a_{1}-3a_{2}+3a_{3}-a_{4}=a_{1}-3a_{1}q+3a_{1}q^{2}-a_{1}q^{3}=a_{1}(1-q)^{3}\).
\item 对任意正整数 \(n\)，有
\[
a_{1}\mathrm C_{n}^{0}-a_{2}\mathrm C_{n}^{1}+a_{3}\mathrm C_{n}^{2}-\cdots+(-1)^{n}a_{n+1}\mathrm C_{n}^{n}=a_{1}(1-q)^{n}.
\]
事实上，等比数列的通项为 \(a_{k+1}=a_{1}q^{k}\)，所以左端等于 \(a_{1}\sum_{k=0}^{n}\mathrm C_{n}^{k}(-q)^{k}\). 根据二项式定理，该式等于 \(a_{1}(1-q)^{n}\)，结论得证.
\end{enumerate}
\end{solution}

\begin{problem}
如图，某隧道设计为双向四车道，车道总宽 \(22\text{ 米}\)，要求通行车辆限高 \(4.5\text{ 米}\)，隧道全长 \(2.5\text{ 千米}\)，隧道的拱线近似地看成半个椭圆形状.

\begin{enumerate}
\item 若最大拱高 \(h\) 为 \(6\text{ 米}\)，则隧道设计的拱宽 \(l\) 是多少？
\item 若最大拱高 \(h\) 不小于 \(6\text{ 米}\)，则应如何设计拱高 \(h\) 和拱宽 \(l\)，才能使半个椭圆形隧道的土方工程量最小？
\end{enumerate}

注：半个椭圆的面积公式为 \(S = \frac{\pi}{4} lh\)，柱体体积为：底面积乘以高，本题结果精确到 \(0.1\text{ 米}\).

\bitmapfigure[max width=0.6\linewidth]{img/2003/shanghai_science/q20_fig1.png}
\end{problem}

\begin{solution}
\begin{enumerate}
\item 取隧道横截面的对称轴为坐标轴，半个椭圆所在椭圆的方程可写为 \(\dfrac{x^{2}}{(l/2)^{2}}+\dfrac{y^{2}}{h^{2}}=1\). 在高度 \(4.5\text{ 米}\) 处，通道宽为 \(22\text{ 米}\)，所以横截面边界点的坐标可取为 \((11,4.5)\)，从而 \(\dfrac{121}{(l/2)^{2}}+\dfrac{4.5^{2}}{h^{2}}=1\). 当 \(h=6\) 时，\(\dfrac{484}{l^{2}}+\dfrac{20.25}{36}=1\)，故 \(l^{2}=\dfrac{7744}{7}\)，即 \(l=\dfrac{88}{\sqrt{7}}\approx33.3\text{ 米}\).
\item 由上面的宽高关系，得 \(\dfrac{484}{l^{2}}+\dfrac{81}{4h^{2}}=1\)，所以 \(l=\dfrac{44h}{\sqrt{4h^{2}-81}}\)，其中 \(h\geqslant6\). 隧道全长固定，土方工程量与半个椭圆的面积 \(S=\dfrac{\pi}{4}lh\) 成正比，因此只需使 \(lh\) 最小. 令 \(t=h^{2}\)，则 \(t\geqslant36\)，且 \(lh=\dfrac{44t}{\sqrt{4t-81}}\). 设 \(F(t)=\dfrac{t}{\sqrt{4t-81}}\)，则
此时 \(F'(t)=\dfrac{2t-81}{(4t-81)^{3/2}}\).
当 \(36\leqslant t<\dfrac{81}{2}\) 时，\(F'(t)<0\)；当 \(t>\dfrac{81}{2}\) 时，\(F'(t)>0\)，故最小值在 \(t=\dfrac{81}{2}\) 处取得. 因而 \(h=\dfrac{9}{\sqrt{2}}\approx6.4\text{ 米}\)，\(l=\dfrac{44}{\sqrt{2}}=22\sqrt{2}\approx31.1\text{ 米}\). 此时半个椭圆的面积为 \(\dfrac{99\pi}{2}\text{ 平方米}\)，相应的最小土方工程量为 \(123750\pi\text{ 立方米}\).
\end{enumerate}
\end{solution}

\begin{problem}
在以 \(O\) 为原点的直角坐标系中，点 \(A(4, -3)\) 为 \(\triangle OAB\) 的直角顶点. 已知 \(|AB| = 2|OA|\)，且点 \(B\) 的纵坐标大于零.

\begin{enumerate}
\item 求向量 \(\overrightarrow{AB}\) 的坐标；
\item 求圆 \(x^{2} - 6x + y^{2} + 2y = 0\) 关于直线 \(OB\) 对称的圆的方程；
\item 是否存在实数 \(a\)，使抛物线 \(y = ax^2 - 1\) 上总有关于直线 \(OB\) 对称的两个点？若不存在，说明理由；若存在，求 \(a\) 的取值范围.
\end{enumerate}
\end{problem}

\begin{solution}
\begin{enumerate}
\item 由 \(A(4,-3)\) 得 \(\overrightarrow{AO}=(-4,3)\)，\(|OA|=5\)，所以 \(|AB|=10\). 设 \(\overrightarrow{AB}=(u,v)\)，因为 \(A\) 是直角顶点，有 \((-4,3)\cdot(u,v)=0\)，即 \(-4u+3v=0\)，又 \(u^{2}+v^{2}=100\). 由 \(v=\dfrac{4u}{3}\) 得 \(\dfrac{25u^{2}}{9}=100\)，所以 \((u,v)=(6,8)\) 或 \((-6,-8)\). 相应的点 \(B\) 的纵坐标分别为 \(5\) 和 \(-11\)，由题意取 \(\overrightarrow{AB}=(6,8)\)，并有 \(B=(10,5)\).
\item 原圆方程可化为 \((x-3)^{2}+(y+1)^{2}=10\)，所以圆心为 \(C_{0}=(3,-1)\)，半径平方为 \(10\). 由 \(B=(10,5)\)，直线 \(OB\) 的方向向量可取 \((2,1)\). 圆心在 \(OB\) 上的射影为 \(H=\dfrac{(3,-1)\cdot(2,1)}{2^{2}+1^{2}}(2,1)=(2,1)\).
关于 \(OB\) 对称后，圆心变为 \(C_{0}'=2H-C_{0}=(1,3)\)，半径不变，故所求圆的方程为 \((x-1)^{2}+(y-3)^{2}=10\)，即 \(x^{2}-2x+y^{2}-6y=0\).
\item 直线 \(OB\) 的方向向量为 \((2,1)\). 设关于该直线对称的两个不同点为 \(P\) 和 \(Q\)，其中点为 \(M=(2t,t)\)，取垂直于 \(OB\) 的方向向量 \((-1,2)\)，则可设 \(P=(2t-r,t+2r)\)，\(Q=(2t+r,t-2r)\)，且 \(r\neq0\). 若两点都在抛物线 \(y=ax^{2}-1\) 上，则
\[
\begin{cases}
t+2r=a(2t-r)^{2}-1,\\
t-2r=a(2t+r)^{2}-1.
\end{cases}
\]
两式相减并利用 \(r\neq0\)，得 \(4r=-8atr\)，所以 \(at=-\dfrac12\). 两式相加后除以 \(2\)，得 \(t=a(4t^{2}+r^{2})-1\). 代入 \(t=-\dfrac{1}{2a}\)，可得 \(ar^{2}=1-\dfrac{3}{2a}\)，即 \(r^{2}=\dfrac{2a-3}{2a^{2}}\). 因为 \(r\neq0\)，必须有 \(r^{2}>0\)，所以 \(a>\dfrac32\). 反之，当 \(a>\dfrac32\) 时，取 \(t=-\dfrac{1}{2a}\)，\(r=\sqrt{\dfrac{2a-3}{2a^{2}}}\)，上述两点均满足抛物线方程，且彼此不同，故存在这样的两个点. 所以 \(a\) 的取值范围为 \(\left(\dfrac32,+\infty\right)\).
\end{enumerate}
\end{solution}

\begin{problem}
已知集合 \(M\) 是满足下列性质的函数 \(f(x)\) 的全体：存在非零常数 \(T\)，对任意 \(x \in \R\)，有 \(f(x + T) = Tf(x)\) 成立.

\begin{enumerate}
\item 函数 \(f(x) = x\) 是否属于集合 \(M\)？说明理由.
\item 设函数 \(f(x) = a^x\)（\(a > 0\)，且 \(a \neq 1\)）的图象与 \(y = x\) 的图象有公共点，证明：\(f(x) = a^x \in M\).
\item 若函数 \(f(x) = \sin kx \in M\)，求实数 \(k\) 的取值范围.
\end{enumerate}
\end{problem}

\begin{solution}
\begin{enumerate}
\item 若 \(f(x)=x\) 属于 \(M\)，则存在非零常数 \(T\)，使得对任意 \(x\) 有 \(x+T=Tx\)，即 \((1-T)x+T=0\). 这要求 \(1-T=0\) 且 \(T=0\)，两者矛盾，所以 \(f(x)=x\) 不属于 \(M\).
\item 设图象有公共点 \((x_{0},x_{0})\)，则 \(a^{x_{0}}=x_{0}\). 由于 \(a^{x_{0}}>0\)，有 \(x_{0}>0\)，可取非零常数 \(T=x_{0}\). 此时 \(a^{T}=T\)，所以对任意 \(x\in\R\)，有 \(f(x+T)=a^{x+T}=a^{x}a^{T}=Tf(x)\)，故 \(f(x)=a^{x}\in M\).
\item 当 \(k=0\) 时，\(f(x)=0\)，任取非零常数 \(T\) 都有 \(f(x+T)=Tf(x)\)，所以 \(k=0\) 满足条件. 设 \(k\neq0\)，则对任意 \(x\) 有 \(\sin(k(x+T))=T\sin(kx)\)，即
\(\sin(kx)\cos(kT)+\cos(kx)\sin(kT)=T\sin(kx)\).
由于 \(\sin(kx)\) 与 \(\cos(kx)\) 的系数必须分别相等，得到 \(\sin(kT)=0\)，\(\cos(kT)=T\). 因而 \(T=1\) 或 \(T=-1\). 当 \(T=1\) 时，\(\cos k=1\)，故 \(k=2n\pi\)；当 \(T=-1\) 时，\(\cos k=-1\)，故 \(k=(2n+1)\pi\)，其中 \(n\in\Z\). 两种情形合并为 \(k=n\pi\)，\(n\in\Z\). 反之，若 \(k=n\pi\)，则当 \(n\) 为偶数取 \(T=1\)，当 \(n\) 为奇数取 \(T=-1\)，均可满足定义中的等式，所以 \(k\) 的取值范围为 \(\pi\Z\).
\end{enumerate}
\end{solution}
